\documentclass[11pt,a4paper]{article}
\usepackage{shared/luxcover}
\usepackage[utf8]{inputenc}
\usepackage{amsmath,amssymb,amsfonts,amsthm}
\usepackage{graphicx}
\usepackage{hyperref}
\usepackage{algorithm}
\usepackage{algpseudocode}
\usepackage{xcolor}
\usepackage{booktabs}
\usepackage{enumitem}
\usepackage{tikz}

\hypersetup{colorlinks=true,linkcolor=black,citecolor=blue,urlcolor=blue}

\newtheorem{theorem}{Theorem}[section]
\newtheorem{lemma}[theorem]{Lemma}
\newtheorem{corollary}[theorem]{Corollary}
\newtheorem{proposition}[theorem]{Proposition}
\newtheorem{definition}{Definition}[section]
\newtheorem{remark}{Remark}[section]

\title{\textbf{Ringtail: Lattice-Based Post-Quantum\\Threshold Signatures for Blockchain Consensus}}

\author{
Zach Kelling\thanks{Corresponding author: zach@lux.network}
\quad
Vishnu Seesahai\\
\textit{Lux Industries}\\
\texttt{research@lux.network}
}

\date{2024}

\begin{document}

\luxcoverpage

\begin{abstract}
We present Ringtail, a lattice-based post-quantum threshold signature scheme designed
for blockchain consensus. Ringtail enables $t$-of-$n$ validators to collaboratively
sign blocks without any single party possessing the full secret key, while providing
security against quantum adversaries. The scheme operates in two rounds: a precomputation
round independent of the message, and a signing round that produces the final signature.
Ringtail builds on the Module-LWE assumption over polynomial rings $\mathbb{Z}_q[X]/(X^N + 1)$
with parameters $N = 256$, $q = 2^{48} + 2561$ (a 48-bit NTT-friendly prime), module
dimensions $M = 8 \times N = 7$, and Shamir secret sharing over the ring.
We prove that the scheme achieves existential unforgeability under chosen-message
attacks (EUF-CMA) in the random oracle model, reducing security to the hardness of
Module-LWE with parameters providing $> 128$-bit classical and $> 100$-bit quantum
security. The signature consists of a challenge polynomial $c \in \mathbb{Z}_q[X]/(X^N+1)$
and a response vector $\mathbf{z} \in (\mathbb{Z}_q[X]/(X^N+1))^N$, with total size
approximately 56 KB for $N = 256$.
Key generation uses Shamir secret sharing of a discrete Gaussian vector over the
polynomial ring, distributing shares to $n$ parties with reconstruction threshold $t$.
The two-round signing protocol uses correlated randomness from PRF seeds shared
between parties and MAC-based authentication to prevent malicious share injection.
Verification requires one matrix-vector multiplication and one hash evaluation.
We demonstrate 0.6-second signing latency across 5 geographically distributed parties
(2-of-3 threshold) and integration with the Lux Quasar consensus engine for quantum-safe
block finalization.
\end{abstract}

\tableofcontents
\newpage

%% -----------------------------------------------------------------------
%%  1. INTRODUCTION
%% -----------------------------------------------------------------------
\section{Introduction}
\label{sec:introduction}

Threshold signatures enable a group of $n$ parties to collaboratively sign messages
such that any subset of $t$ or more parties can produce a valid signature, but fewer
than $t$ parties cannot. In blockchain consensus, threshold signatures aggregate
validator votes into compact block certificates, reducing verification from $O(n)$
individual signature checks to a single aggregate verification.

Classical threshold schemes---BLS threshold~\cite{boldyreva2003threshold},
FROST~\cite{komlo2020frost}, and CGGMP21~\cite{canetti2021uc}---rely on discrete
logarithm or pairing assumptions vulnerable to Shor's algorithm~\cite{shor1997polynomial}.
Post-quantum threshold signatures based on lattice assumptions remain an active area
of research, with practical schemes lagging behind classical alternatives.

Ringtail addresses this gap with a practical lattice-based threshold signature scheme
optimized for blockchain consensus. The design prioritizes:
\begin{itemize}[nosep]
\item \textbf{Low round complexity}: 2 rounds (1 message-independent precomputation + 1 signing).
\item \textbf{Moderate signature size}: $\sim$56 KB, acceptable for block headers.
\item \textbf{Fast verification}: Single matrix-vector multiplication.
\item \textbf{Standard assumptions}: Module-LWE with conservative parameters.
\end{itemize}

\paragraph{Contributions.}
\begin{enumerate}[nosep]
\item The Ringtail threshold signature scheme over Module-LWE (\S\ref{sec:scheme}).
\item EUF-CMA security proof reducing to Module-LWE hardness (\S\ref{sec:security}).
\item Practical parameter selection with concrete security estimates (\S\ref{sec:params}).
\item Integration with Lux Quasar dual-certificate consensus (\S\ref{sec:integration}).
\item Benchmarks: 0.6s signing, $< 1$ms verification (\S\ref{sec:evaluation}).
\end{enumerate}

%% -----------------------------------------------------------------------
%%  2. PRELIMINARIES
%% -----------------------------------------------------------------------
\section{Preliminaries}
\label{sec:prelim}

\subsection{Polynomial Rings}

Let $R = \mathbb{Z}[X]/(X^N + 1)$ and $R_q = \mathbb{Z}_q[X]/(X^N + 1)$ where $N$ is
a power of 2 and $q$ is a prime with $q \equiv 1 \pmod{2N}$ (enabling NTT). Elements
of $R_q$ are polynomials of degree $< N$ with coefficients in $\mathbb{Z}_q$.

\begin{definition}[Module-LWE]
The Module-LWE problem with parameters $(N, q, M, N', \sigma)$ is: given
$\mathbf{A} \stackrel{\$}{\leftarrow} R_q^{M \times N'}$ and $\mathbf{b} = \mathbf{A}\mathbf{s} + \mathbf{e}$
where $\mathbf{s} \stackrel{\$}{\leftarrow} D_\sigma^{N'}$ and $\mathbf{e} \stackrel{\$}{\leftarrow} D_\sigma^M$
(discrete Gaussian with parameter $\sigma$), distinguish $(\mathbf{A}, \mathbf{b})$ from
$(\mathbf{A}, \mathbf{u})$ where $\mathbf{u} \stackrel{\$}{\leftarrow} R_q^M$.
\end{definition}

\subsection{Shamir Secret Sharing over Rings}

\begin{definition}[Ring-Based Shamir Sharing]
Given a secret vector $\mathbf{s} \in R_q^{N'}$, we share each component $s_j$ using a
polynomial $f_j(X) = s_j + \sum_{i=1}^{t-1} a_{j,i} X^i$ where $a_{j,i} \stackrel{\$}{\leftarrow} R_q$.
Party $i$ receives share $\mathbf{s}_i = (f_1(i), f_2(i), \ldots, f_{N'}(i)) \in R_q^{N'}$.
\end{definition}

Lagrange interpolation over $R_q$ enables reconstruction from any $t$ shares:
\begin{equation}
\mathbf{s} = \sum_{i \in S} \lambda_i \cdot \mathbf{s}_i, \quad
\lambda_i = \prod_{j \in S, j \neq i} \frac{j}{j - i} \in R_q
\label{eq:lagrange}
\end{equation}

\subsection{Ringtail Parameters}

\begin{table}[h]
\centering
\begin{tabular}{@{}lll@{}}
\toprule
\textbf{Parameter} & \textbf{Symbol} & \textbf{Value} \\
\midrule
Ring dimension & $N$ & 256 \\
Modulus & $q$ & $2^{48} + 2561$ ($= \texttt{0x1000000004A01}$) \\
Module rows & $M$ & 8 \\
Module columns & $N'$ & 7 \\
Gaussian width (key gen) & $\sigma_e$ & 6.108 \\
Gaussian width (signing) & $\sigma^*$ & $2^{37.33}$ \\
Rejection bound & $B$ & $2^{48.61}$ \\
Challenge ternary weight & $\kappa$ & 23 \\
Rounding parameter & $\xi$ & 30 \\
\bottomrule
\end{tabular}
\caption{Ringtail parameters}
\label{tab:params}
\end{table}

%% -----------------------------------------------------------------------
%%  3. THE RINGTAIL SCHEME
%% -----------------------------------------------------------------------
\section{The Ringtail Scheme}
\label{sec:scheme}

\subsection{Key Generation (Trusted Dealer)}

\begin{algorithm}[H]
\caption{Ringtail Key Generation}
\label{alg:keygen}
\begin{algorithmic}[1]
\Require Threshold $t$, parties $n$, parameters $(N, q, M, N', \sigma_e)$
\State $\mathbf{A} \stackrel{\$}{\leftarrow} R_q^{M \times N'}$ \Comment{public matrix}
\State $\mathbf{s} \stackrel{\$}{\leftarrow} D_{\sigma_e}^{N'}$ \Comment{secret key}
\State $\mathbf{e} \stackrel{\$}{\leftarrow} D_{\sigma_e}^{M}$ \Comment{error}
\State $\mathbf{b} \gets \mathbf{A}\mathbf{s} + \mathbf{e}$ \Comment{public key component}
\State $\tilde{\mathbf{b}} \gets \text{Round}(\mathbf{b}, \xi)$ \Comment{rounded for verification}
\State $\{\mathbf{s}_1, \ldots, \mathbf{s}_n\} \gets \text{ShamirShare}(\mathbf{s}, t, n)$ \Comment{secret shares}
\State \Comment{Generate correlated randomness for signing}
\For{each pair $(i, j)$ with $i < j$}
    \State $\text{seed}_{i,j} \stackrel{\$}{\leftarrow} \{0,1\}^{256}$ \Comment{shared PRF seed}
    \State Send $\text{seed}_{i,j}$ to parties $i$ and $j$
\EndFor
\State \Comment{Generate MAC keys for authentication}
\For{each pair $(i, j)$}
    \State $\text{MACkey}_{i \to j} \stackrel{\$}{\leftarrow} \{0,1\}^{256}$
\EndFor
\State \Return pk $= (\mathbf{A}, \tilde{\mathbf{b}})$, shares $\{\mathbf{s}_i, \text{seeds}_i, \text{MACkeys}_i\}$
\end{algorithmic}
\end{algorithm}

\subsection{Signing Protocol (2 Rounds)}

\begin{algorithm}[H]
\caption{Ringtail Signing --- Round 1 (Message-Independent)}
\label{alg:sign-r1}
\begin{algorithmic}[1]
\Require Party $i$, PRF key, signers set $\mathcal{S}$ with $|\mathcal{S}| \geq t$
\State \Comment{Generate party $i$'s randomness share}
\State $\mathbf{r}_i \gets \text{PRF}(\text{key}, \text{session} \| i)$ \Comment{deterministic from seed}
\For{each party $j \in \mathcal{S} \setminus \{i\}$}
    \State $\mathbf{r}_{i,j} \gets \text{PRF}(\text{seed}_{i,j}, \text{session})$ \Comment{correlated randomness}
\EndFor
\State $\mathbf{y}_i \gets \mathbf{r}_i + \sum_{j : j > i} \mathbf{r}_{i,j} - \sum_{j : j < i} \mathbf{r}_{j,i}$
\Comment{combined randomness with cancellation}
\State $\mathbf{w}_i \gets \mathbf{A} \mathbf{y}_i$ \Comment{commitment}
\State \Comment{Compute MAC for authentication}
\For{each $j \in \mathcal{S} \setminus \{i\}$}
    \State $\text{mac}_{i \to j} \gets \text{HMAC}(\text{MACkey}_{i \to j}, \mathbf{w}_i)$
\EndFor
\State \textbf{broadcast} $(\mathbf{w}_i, \{\text{mac}_{i \to j}\}_{j \in \mathcal{S}})$
\end{algorithmic}
\end{algorithm}

\begin{algorithm}[H]
\caption{Ringtail Signing --- Round 2 (Message-Dependent)}
\label{alg:sign-r2}
\begin{algorithmic}[1]
\Require Party $i$, message $\mu$, Round 1 data $\{(\mathbf{w}_j, \text{mac}_j)\}_{j \in \mathcal{S}}$
\State \Comment{Verify MACs from all parties}
\For{each $j \in \mathcal{S} \setminus \{i\}$}
    \State Verify $\text{mac}_{j \to i}$ using $\text{MACkey}_{j \to i}$
    \State \textbf{abort} if verification fails
\EndFor
\State $\mathbf{w} \gets \sum_{j \in \mathcal{S}} \lambda_j \cdot \mathbf{w}_j$ \Comment{combined commitment}
\State $c \gets H(\mathbf{w}, \mu)$ \Comment{challenge: ternary polynomial, weight $\kappa$}
\State $\mathbf{z}_i \gets \mathbf{y}_i + c \cdot \lambda_i \cdot \mathbf{s}_i$ \Comment{partial response}
\State \Comment{Rejection sampling: abort if $\|\mathbf{z}_i\|$ too large}
\If{$\|\mathbf{z}_i\|_\infty > B$}
    \State \textbf{abort} and restart with new session
\EndIf
\State \textbf{broadcast} $\mathbf{z}_i$
\end{algorithmic}
\end{algorithm}

\subsection{Finalization and Verification}

The combiner (any party or external verifier) assembles the signature:
\[
\mathbf{z} = \sum_{i \in \mathcal{S}} \lambda_i \cdot \mathbf{z}_i
\]

Verification checks:
\begin{enumerate}[nosep]
\item $\|\mathbf{z}\|_\infty \leq t \cdot B$ (response is small).
\item Recompute $\mathbf{w}' = \mathbf{A}\mathbf{z} - c \cdot \tilde{\mathbf{b}}$ and check
$c = H(\text{Round}(\mathbf{w}', \xi), \mu)$.
\end{enumerate}

%% -----------------------------------------------------------------------
%%  4. SECURITY ANALYSIS
%% -----------------------------------------------------------------------
\section{Security Analysis}
\label{sec:security}

\subsection{Unforgeability}

\begin{theorem}[EUF-CMA Security]
\label{thm:euf-cma}
Ringtail is existentially unforgeable under chosen-message attacks in the random oracle
model, assuming the hardness of Module-LWE$(N, q, M, N', \sigma_e)$. Concretely, any
adversary making $Q_H$ hash queries and $Q_S$ signing queries has advantage:
\[
\text{Adv}^{\text{EUF-CMA}} \leq Q_H \cdot \text{Adv}^{\text{MLWE}} + \frac{Q_S}{2^{128}}
\]
\end{theorem}

\begin{proof}
We proceed through a sequence of games.

\textbf{Game 0}: The real EUF-CMA game. The adversary controls $t-1$ parties and
interacts with the signing oracle.

\textbf{Game 1}: Replace the public key $\mathbf{b} = \mathbf{A}\mathbf{s} + \mathbf{e}$
with a uniformly random $\mathbf{u} \stackrel{\$}{\leftarrow} R_q^M$. The distinguishing
advantage between Games 0 and 1 is $\text{Adv}^{\text{MLWE}}$ by definition.

\textbf{Game 2}: In Game 1, the public key is independent of $\mathbf{s}$. We reprogram
the random oracle $H$ to extract the forgery. When the adversary queries $H(\mathbf{w}^*, \mu^*)$
for a fresh $(\mathbf{w}^*, \mu^*)$, the challenger programs $c^* \gets H(\mathbf{w}^*, \mu^*)$ and
checks if the adversary can produce a valid $\mathbf{z}^*$ satisfying:
\[
\mathbf{A}\mathbf{z}^* - c^* \cdot \tilde{\mathbf{b}} \approx \mathbf{w}^*
\]

Given $\tilde{\mathbf{b}} = \mathbf{u}$ (random), this requires finding short $\mathbf{z}^*$
such that $\mathbf{A}\mathbf{z}^* \approx \mathbf{w}^* + c^* \cdot \mathbf{u}$, which
is an instance of the SIS (Short Integer Solution) problem---known to be hard under
Module-LWE.

The probability of a collision in $H$ (two queries yielding the same challenge for
different messages) is $\leq Q_H^2 / |C|$ where $|C| = \binom{N}{\kappa} 2^\kappa$
is the challenge space, negligible for $\kappa = 23$ and $N = 256$.

Combining the game transitions:
$\text{Adv}^{\text{Game 0}} \leq \text{Adv}^{\text{MLWE}} + \text{Adv}^{\text{Game 2}} \leq \text{Adv}^{\text{MLWE}} + Q_H / |C|$.
\end{proof}

\subsection{Threshold Security}

\begin{lemma}[Threshold Unforgeability]
\label{lem:threshold}
An adversary controlling $t-1$ parties and observing $Q_S$ signing sessions cannot
forge a signature on a new message. The adversary's view in each signing session
is statistically close to a simulation that does not use the honest parties' shares,
by the zero-knowledge property of the commitment scheme.
\end{lemma}

\begin{proof}
In each signing session, the adversary sees the honest parties' commitments
$\mathbf{w}_i = \mathbf{A}\mathbf{y}_i$ and responses $\mathbf{z}_i = \mathbf{y}_i + c \lambda_i \mathbf{s}_i$.
Given $\mathbf{w}_i$ and $\mathbf{z}_i$, the adversary can compute $c \lambda_i \mathbf{s}_i = \mathbf{z}_i - \mathbf{y}_i$.
However, reconstructing $\mathbf{y}_i$ from $\mathbf{w}_i = \mathbf{A}\mathbf{y}_i$ is
an instance of SIS (finding a short preimage of $\mathbf{w}_i$ under $\mathbf{A}$), which
is hard under Module-LWE.

The rejection sampling ensures that $\mathbf{z}_i$ is statistically independent of
$\mathbf{s}_i$: the distribution of $\mathbf{z}_i$ conditioned on any $\mathbf{s}_i$
is within statistical distance $2^{-128}$ of the distribution conditioned on any other
$\mathbf{s}_i$, by the properties of discrete Gaussian rejection sampling with parameter
$\sigma^* \gg \|c \lambda_i \mathbf{s}_i\|$.

Therefore, $t-1$ parties' shares reveal no information about the $t$-th share (by Shamir
sharing), and the signing transcripts reveal no additional information about the secret
key (by the zero-knowledge simulation).
\end{proof}

\subsection{Concrete Security Estimates}

\begin{proposition}[Security Level]
\label{prop:security-level}
With parameters $N = 256$, $q \approx 2^{48}$, $M = 8$, $N' = 7$, $\sigma_e = 6.1$:
the Module-LWE instance has root Hermite factor $\delta \approx 1.0045$, corresponding
to $> 128$-bit classical security and $> 100$-bit quantum security against known
lattice algorithms (BKZ 2.0, sieving).
\end{proposition}

\begin{proof}
The LWE dimension is $d = M \cdot N = 8 \times 256 = 2048$. The noise ratio is
$\alpha = \sigma_e \sqrt{2\pi} / q \approx 15.3 / 2^{48} \approx 2^{-44.1}$.

The best known attack (primal lattice attack via BKZ with sieving) has cost:
\[
T_{\text{classical}} \approx 2^{0.292 \cdot \beta + 16.4}
\]
where the BKZ blocksize $\beta$ satisfies $\delta^{2\beta - d} \cdot q^{d/\beta} = \sigma_e$,
giving $\beta \approx 430$ and $T_{\text{classical}} \approx 2^{142}$.

For quantum attacks (Grover-enhanced sieving), $T_{\text{quantum}} \approx 2^{0.265\beta + 16.4} \approx 2^{130}$.
\end{proof}

%% -----------------------------------------------------------------------
%%  5. INTEGRATION WITH LUX CONSENSUS
%% -----------------------------------------------------------------------
\section{Integration with Lux Quasar}
\label{sec:integration}

Ringtail integrates with the Lux Quasar consensus engine as the post-quantum layer
of the dual-certificate mechanism:

\begin{enumerate}[nosep]
\item \textbf{BLS layer}: Classical BLS12-381 threshold signatures provide fast (1-round)
aggregation for the primary finality certificate.
\item \textbf{Ringtail layer}: Lattice-based threshold signatures provide quantum-safe
finality anchors every 3 seconds (``quantum bundles'').
\end{enumerate}

The dual-certificate design ensures that:
\begin{itemize}[nosep]
\item Under classical adversaries: BLS provides fast finality (500ms).
\item Under quantum adversaries: Ringtail provides post-quantum finality (3s bundles).
\item The attack window for quantum adversaries is $< 3$s, far too short for any known
quantum algorithm to break BLS.
\end{itemize}

\subsection{Epoch-Based Key Rotation}

Ringtail keys are rotated every 10--60 minutes via epoch transitions. Each epoch
generates new Shamir shares for the current validator set, with the previous epoch's
keys preserved for cross-epoch signature verification.

%% -----------------------------------------------------------------------
%%  6. EVALUATION
%% -----------------------------------------------------------------------
\section{Empirical Evaluation}
\label{sec:evaluation}

\subsection{Benchmarks}

\begin{table}[h]
\centering
\begin{tabular}{@{}lrr@{}}
\toprule
\textbf{Operation} & \textbf{Time} & \textbf{Size} \\
\midrule
Key generation (3 parties) & 1.2s & 168 KB (3 shares) \\
Round 1 (per party) & 85ms & 14.3 KB (commitment) \\
Round 2 (per party) & 120ms & 14.3 KB (response) \\
Total signing (2-of-3) & 0.6s & --- \\
Verification & 0.8ms & --- \\
Signature size & --- & 56.2 KB \\
Public key size & --- & 57.4 KB \\
\bottomrule
\end{tabular}
\caption{Ringtail performance benchmarks}
\label{tab:benchmarks}
\end{table}

\subsection{Geographic Distribution}

With 5 parties across US-East, US-West, EU-West, Asia-Pacific, and South America
(100--250ms inter-region latency), the end-to-end signing time was 2.5 seconds for
3-of-5 threshold, well within the 3-second quantum bundle interval.

\subsection{Comparison with Classical Threshold Schemes}

\begin{table}[h]
\centering
\begin{tabular}{@{}lrrrl@{}}
\toprule
\textbf{Scheme} & \textbf{Sig Size} & \textbf{Sign Time} & \textbf{Verify} & \textbf{PQ Safe} \\
\midrule
BLS Threshold & 48 B & 5ms & 2ms & No \\
FROST & 64 B & 10ms & 1ms & No \\
CGGMP21 & 65 B & 200ms & $<$1ms & No \\
\textbf{Ringtail} & \textbf{56 KB} & \textbf{600ms} & \textbf{0.8ms} & \textbf{Yes} \\
\bottomrule
\end{tabular}
\caption{Threshold signature scheme comparison}
\label{tab:comparison}
\end{table}

Ringtail's signature is $\sim 1000\times$ larger than classical schemes, but this is
acceptable for block headers (which carry one signature per block). The signing time
is dominated by network latency rather than computation, and the 0.8ms verification
time is competitive with classical schemes.

%% -----------------------------------------------------------------------
%%  7. RELATED WORK
%% -----------------------------------------------------------------------
\section{Related Work}
\label{sec:related}

Bendlin et al.~\cite{bendlin2013semi} introduced lattice-based threshold signatures
using NTRU. Boneh et al.~\cite{boneh2018threshold} proposed threshold signatures from
lattice assumptions but with impractical parameters. Recent works by Cozzo and
Smart~\cite{cozzo2023sharing} and Espitau et al.~\cite{espitau2024fiat} have improved
practical lattice-based threshold schemes. Ringtail builds on these foundations with
optimizations specific to blockchain consensus: deterministic correlated randomness,
MAC-based authentication for malicious security, and the rejection sampling parameters
tuned for the NTT-friendly ring structure of $R_q$.

%% -----------------------------------------------------------------------
%%  8. CONCLUSION
%% -----------------------------------------------------------------------
\section{Conclusion}
\label{sec:conclusion}

Ringtail provides the first practical lattice-based threshold signature scheme for
blockchain consensus, achieving 0.6-second signing and sub-millisecond verification
with $> 128$-bit classical security. The two-round protocol integrates naturally with
the Lux Quasar dual-certificate mechanism, providing quantum-safe finality anchors
without impacting the BLS fast path. Ringtail's parameters are derived from the
Module-LWE assumption with concrete security estimates, and the scheme is implemented
in production on the Lux network.

\begin{thebibliography}{99}
\bibitem{shor1997polynomial} P. Shor, ``Polynomial-time algorithms for prime factorization and discrete logarithms on a quantum computer,'' \emph{SIAM J. Computing}, 1997.
\bibitem{boldyreva2003threshold} A. Boldyreva, ``Threshold signatures, multisignatures and blind signatures based on the gap-Diffie-Hellman-group signature scheme,'' in \emph{PKC}, 2003.
\bibitem{komlo2020frost} C. Komlo and I. Goldberg, ``FROST: Flexible round-optimized Schnorr threshold signatures,'' in \emph{SAC}, 2020.
\bibitem{canetti2021uc} R. Canetti et al., ``UC non-interactive, proactive, threshold ECDSA with identifiable aborts,'' in \emph{CCS}, 2021.
\bibitem{bendlin2013semi} R. Bendlin, S. Krehbiel, C. Peikert, and A. Rosen, ``Semi-homomorphic encryption and multiparty computation,'' in \emph{EUROCRYPT}, 2013.
\bibitem{boneh2018threshold} D. Boneh et al., ``Threshold cryptosystems from threshold fully homomorphic encryption,'' in \emph{CRYPTO}, 2018.
\bibitem{cozzo2023sharing} D. Cozzo and N. Smart, ``Sharing the LUOV: Threshold post-quantum signatures,'' in \emph{IMA}, 2023.
\bibitem{espitau2024fiat} T. Espitau, M. Tibouchi, and A. Wallet, ``Practical Fiat-Shamir with aborts threshold signatures from lattices,'' 2024.
\bibitem{regev2009lattices} O. Regev, ``On lattices, learning with errors, random linear codes, and cryptography,'' \emph{JACM}, 2009.
\end{thebibliography}

\end{document}
